\subsection{积分表示的唯一性}

设 E 是一个 Hausdorff 弱局部凸空间（TVS, II, §6, No.~2），C 是 E 中一个真尖凸锥。已知 C 是 E 中某个与 E 的向量空间结构相容的序关系下的 \(\geqslant0\) 元素组成之集。当说 C 是格序的时候，当然理解为 E 的序在 C 上诱导的序。

引理 5. --- 假定 C 是弱完备的。设 A 是 E 上连续线性形式在 C 上的限制所成之集。设 \((f_{\lambda})_{\lambda \in \Lambda}\) 是 A 中元素的一个有限族，且 \(f = \sup(f_{\lambda})\)。对每个 \(x \in C\)，定义

\[
\overline {{{f}}} (x) = \sup \left(f (x _ {1}) + f (x _ {2}) + \dots + f (x _ {n})\right),
\]

其中上确界取遍集 \(S_{x}\)——由 C 中元素的序列 \((x_{1}, x_{2}, \ldots, x_{n})\)（满足 \(x_{1} + x_{2} + \cdots + x_{n} = x\)）组成。设 \(\operatorname{Card}\Lambda = p\)。那么存在 \((y_{1}, \ldots, y_{p}) \in \operatorname{S}_{x}\) 使得 \(\overline{f}(x) = f(y_{1}) + \cdots + f(y_{p})\)。

用 \(f_1, f_2, \ldots, f_p\) 表示族 \((f_\lambda)\) 的元素。对 \(k = 1, 2, \ldots, p\)，设 \(C_k\) 是由满足下列关系的 \(y \in C\) 组成之集：

\[
f _ {1} (y) <   f (y), f _ {2} (y) <   f (y), \dots , f _ {k - 1} (y) <   f (y), f _ {k} (y) = f (y).
\]

诸 \(C_{k}\) 是两两不相交的凸锥，其并为 C。设 C 中的 \(x_{1}, x_{2}, \ldots, x_{n}\) 满足 \(x_{1} + x_{2} + \cdots + x_{n} = x\)。设 \(y_{k}\) 是那些 \(x_{i}\)（属于 \(C_{k}\) 的）之和。那么 \(y_{1} + y_{2} + \cdots + y_{p} = x\)。由于 f 在 \(C_{k}\) 上是仿射的，故 \(f(y_{1}) + \cdots + f(y_{p}) = f(x_{1}) + \cdots + f(x_{n})\)。因此

\[
f (x) = \sup \bigl (f (y _ {1}) + \dots + f (y _ {p}) \bigr),\tag{19}
\]

其中 \((y_{1}, y_{2}, \ldots, y_{p})\) 取遍由满足 \(y_{1} + y_{2} + \cdots + y_{p} = x\) 的 C 中 p 个点的序列组成之集。令 \(D = C \cap (x - C)\)。由于 D 是紧的（TVS, II, §6, No.~8, 命题 11 的推论 2），元素 \((y_{1}, \ldots, y_{p})\)（\(D^{p}\) 中的，满足 \(y_{1} + \cdots + y_{p} = x\)）所成之集也是紧的，因而上确界 (19) 可以达到。

引理 6. --- 我们保持引理 5 的假设与记号，并假定诸 \(f_{\lambda}\) 是正的。函数 \(\overline{f}\) 在 C 中是正齐次的、凹的且上半连续的。若 C 是格序的，则它是仿射的。

显然 \(\overline{f}\) 是正齐次的。设 x, y 属于 C。若 C 中的 \(x_{1}, \ldots, x_{m}, y_{1}, \ldots, y_{n}\) 满足 \(x_{1} + \cdots + x_{m} = x\)，\(y_{1} + \cdots + y_{n} = y\)，则 \(x_{1} + \cdots + x_{m} + y_{1} + \cdots + y_{n} = x + y\)，因此

\[
f (x _ {1}) + \dots + f (x _ {m}) + f (y _ {1}) + \dots + f (y _ {n}) \leqslant \overline {{f}} (x + y);
\]

由此得出 \(\overline{f}(x) + \overline{f}(y) \leqslant \overline{f}(x + y)\)，因此 \(\overline{f}\) 是凹的。设 \(L\)（相应地，\(L_{\lambda}\)）是由 \((t, x) \in \mathbf{R} \times E\)（满足 \(x \in C\) 且 \(0 \leqslant t \leqslant \overline{f}(x)\)（相应地，\(0 \leqslant t \leqslant f_{\lambda}(x)\)））组成之集。每个 \(L_{\lambda}\) 都在弱完备真凸锥 \(\mathbf{R}_{+} \times C\) 中是闭的，因此和 \(\sum_{\lambda \in \Lambda} L_{\lambda}\) 是闭的（TVS, II, §6, No.~8, 命题 11 的推论 2）。由引理 5，这个和等于 L。于是 L 是闭的，这就证明 \(\overline{f}\) 是上半连续的。最后，假定 C 是格序的，我们来证明 \(\overline{f}\) 是凸的。设 \(x, y\) 属于 C 且 \(\varepsilon > 0\)。存在 C 中的 \(z_1, z_2, \ldots, z_n\)，使得 \(f(z_1) + \cdots + f(z_n) \geqslant \overline{f}(x + y) - \varepsilon\) 且 \(z_1 + \cdots + z_n = x + y\)。向量空间 C - C 对 E 的序所诱导的序是格序的（A, VI, §1, No.~9, 命题 8）。由分解定理（同前，No.~10, 定理 1），存在 C 中的 \(x_1, \ldots, x_n, y_1, \ldots, y_n\)，使得

\[
x _ {1} + y _ {1} = z _ {1}, \dots , x _ {n} + y _ {n} = z _ {n}, x _ {1} + \dots + x _ {n} = x, y _ {1} + \dots + y _ {n} = y.
\]

于是，由于 \(f\) 是正齐次且凸的，

\[
\begin{array}{r l} \overline {{f}} (x + y) & \leqslant \varepsilon + f (z _ {1}) + \dots + f (z _ {n}) \\ & \leqslant \varepsilon + f (x _ {1}) + f (y _ {1}) + \dots + f (x _ {n}) + f (y _ {n}) \\ & \leqslant \varepsilon + \overline {{f}} (x) + \overline {{f}} (y). \end{array}
\]

由于 \(\varepsilon\) 是任意的数 \(>0\)，我们已证明 \(\overline{f}\) 的确是凸的。

定理 3 (Choquet). --- 设 E 是一个 Hausdorff 弱局部凸空间，C 是 E 中以 0 为顶点的弱完备真凸锥，G 是 C 的极值母线之并，K 是 C 的一个紧凸子集，\(\lambda\) 与 \(\lambda'\) 是 K 上质量为 1、具有相同重心的正测度，且满足 \(\lambda^{*}(\mathrm{K} - (\mathrm{K} \cap \mathrm{G})) = \lambda'^{*}(\mathrm{K} - (\mathrm{K} \cap \mathrm{G})) = 0\)。假定 C 是格序的。那么，对 C 上每个下半连续、正齐次的凸函数 \(f \geqslant 0\)，有 \(\lambda^{*}(f|K) = \lambda'^{*}(f|K)\)。

设 \(\mathcal{A}\)（相应地，\(\mathcal{A}'\)）是 E 上连续线性形式（相应地，仿射函数）在 \(\dot{\mathrm{C}}\) 上的限制所成之集。我们知道（TVS, II, §5, No.~4, 注 2）\(f\) 是 \(\mathcal{A}\) 中 \(\leqslant f\) 的元素所成之集的上包络。形如 \(\sup(f_1, \ldots, f_p)\) 的函数所成之集（其中 \(f_1, \ldots, f_p\) 属于 \(\mathscr{A}\)，\(f_1 \geqslant 0, \ldots, f_p \geqslant 0\)）是一个递增有向集，且以 \(f\) 为其上包络。考虑到 §1, No.~1, 定理 1，只须验证等式 \(\lambda(f|K) = \lambda'(f|K)\)（当 \(f\) 具有前述形式时）。

按引理 5 那样定义 \(\overline{f}\)。显然，\(\overline{f}(y) = f(y)\)（若 \(y \in G\)）。由于 \(\lambda^{*}(\mathrm{K} - (\mathrm{K} \cap \mathrm{G})) = 0\)，我们有 \(\lambda(f|K) = \lambda(\overline{f}|K)\)。由引理 6，\(\overline{f}\) 是仿射的且上半连续。因此 \(\overline{f}|K\) 是 \(\mathcal{A}'\) 的元素在 K 上的限制所成的一个递减有向集的下包络（TVS, II, §5, No.~4, 命题 6）。设 \(x \in K\) 是 \(\lambda\) 的重心。若 \(g \in \mathcal{A}\)，则 \(\lambda(g|K) = g(x)\)。因此 \(\lambda(\overline{f}|K) = \overline{f}(x)\)（§4, No.~4, 命题 5 的推论 2）。于是 \(\lambda(f|K) = \overline{f}(x)\)，并且同理可见 \(\lambda'(f|K) = \overline{f}(x)\)。

推论. --- 设 E 是一个 Hausdorff 局部凸空间，C 是 E 中以 0 为顶点、具有紧底 M 的真凸锥，并设 G 是 C 的极值母线之并。设 \(x \in M\)。若 C 是格序的，则 M 上至多存在一个质量为 1 的正测度 \(\lambda\)，使得 \(\lambda^{*}(\mathrm{M} - (\mathrm{G} \cap \mathrm{M})) = 0\)，且以 x 为重心。

把 E 的拓扑换成弱化的拓扑（这不改变 M 的拓扑），就可以假定 E 是弱空间。设 \(\lambda\) 与 \(\lambda'\) 是 M 上具有所述性质的两个测度，并设 h 是 E 上使 M 为 C 与方程 \(h(x)=1\) 的超平面之交的连续线性形式。设 S 是 \(\mathcal{C}(\mathrm{M})\) 中由 C 上正齐次且连续的 \(\geqslant0\) 凸函数在 M 上的限制组成的子集。锥 C 是弱完备的（TVS, II, §7, No.~3）。由定理 3，\(\lambda(f)=\lambda'(f)\)（对每个 \(f\in S\)）。

若 \(f_{1}, f_{2}, f_{3}, f_{4}\) 属于 \(\mathcal{S}\)，则

\[
\sup \left(f _ {1} - f _ {2}, f _ {3} - f _ {4}\right) = \sup \left(f _ {1} + f _ {4}, f _ {3} + f _ {2}\right) - \left(f _ {2} + f _ {4}\right) \in \mathscr {S} - \mathscr {S}
\]

\[
\inf \left(f _ {1} - f _ {2}, f _ {3} - f _ {4}\right) = - \sup \left(f _ {2} - f _ {1}, f _ {4} - f _ {3}\right) \in \mathscr {S} - \mathscr {S}.
\]

由于 \(h|M \in \mathcal{S}\)，\(\mathcal{S} - \mathcal{S}\) 包含常值函数。若 \(x\) 与 \(y\) 是 \(M\) 的两个不同点，则存在 \(E\) 上的一个连续线性形式，它分离 \(x\) 与 \(y\)，且这一形式是两个在 \(C\) 上为正的连续线性形式之差（TVS, II, §6, No.~8, 引理 1）。由上述事实推出：对实数 \(\alpha, \beta\)，存在 \(f \in \mathcal{S} - \mathcal{S}\) 使得 \(f(x) = \alpha\)，\(f(y) = \beta\)。于是 \(\mathcal{S} - \mathcal{S}\) 在 \(\mathcal{C}(M)\) 中对一致收敛拓扑是稠密的（GT, X, §4, No.~1, 命题 2 的推论）。由于 \(\lambda\) 与 \(\lambda'\) 在 \(\mathcal{S} - \mathcal{S}\) 上一致，我们有 \(\lambda = \lambda'\)。

